\section{Stream Processing and Integration}

Batch processing is useful for historical data, but many ML systems also need
fresh features from recent events. The stream layer answers the question: how
can new events update derived data incrementally?

MLStore-Lite implements a small append-only event log, producers, consumers, an
offset store, tumbling windows, and a stream feature processor. Producers append
events to the log. Consumers read from the log and track their position using
offsets. The processor groups events into fixed-size windows and writes updated
features into the same sharded replicated store used by the batch layer.

This is similar in spirit to Kafka and stream processors such as Flink or Kafka
Streams. Kafka provides the durable distributed event log, while Flink and Kafka
Streams provide continuous computations over event streams. MLStore-Lite keeps a
small local version of the same idea: new events are appended, consumed, grouped,
and converted into feature updates.

\subsection{Integrated System}

The integration layer connects all previous components into one object:

\begin{verbatim}
MLStoreLiteSystem
\end{verbatim}

This object creates the sharded replicated store, the batch feature job, the
event log, producer, consumer, offset store, and stream processor. Its purpose is
not to hide the architecture, but to make the whole prototype runnable without
manually wiring every class each time.

The integrated system supports the following workflow:

\begin{enumerate}
    \item compute batch features from historical events;
    \item produce new events into the stream log;
    \item process stream events into windowed features;
    \item read features from the sharded replicated store;
    \item inspect shard distribution and replica status.
\end{enumerate}

This gives the project an end-to-end ML infrastructure story: raw events become
features, and the features are stored in a distributed-system-inspired storage
backend.

\section{Evaluation and Observability}

The evaluation layer adds local observability. In this project, observability
means structured logs, timing measurements, and reproducible experiment records.
The goal is not to build a monitoring platform. The goal is to make the
prototype inspectable.

The Week 8 evaluation script builds the integrated system and measures:

\begin{itemize}
    \item batch feature computation time;
    \item stream event production time;
    \item stream processing time;
    \item feature read latency;
    \item key distribution across shards.
\end{itemize}

Each measurement is written to a JSON-lines file with an experiment name,
metric, value, unit, timestamp, and parameters. The parameters include the
number of events, number of shards, and replication factor. This is important
for reproducibility: a runtime number is more meaningful when the conditions
that produced it are also recorded.

The evaluation script also emits structured JSON logs. A structured log is more
useful than plain text because it can be parsed by other tools. For example, the
system can log that a measurement was recorded together with the operation name,
elapsed time, and output count.

\section{Online Feature Serving and Model Inference}

The final extension adds a small online inference layer. This layer uses the
features computed by the batch and stream processors and serves them to a
deterministic purchase-intent model.

The feature server reads expected user features from MLStore-Lite, including
historical batch features and recent windowed stream features. It also records
whether expected inputs are missing. The model uses the feature values to
produce a purchase probability, but also reports a confidence value and warnings
when serving-time inputs are unavailable.

This is intentionally not a large model-training system. The purpose is to show
the role of the data infrastructure in model serving:

\begin{verbatim}
events -> features -> feature store -> online inference -> prediction log
\end{verbatim}

Predictions are written to a JSON-lines log with the user id, model version,
input features, prediction probability, confidence, label, and warnings. This
connects the AI extension back to MLOps: model outputs should be traceable and
debuggable, not just returned and forgotten.

\section{Comparison With Production Tools}

MLStore-Lite is a teaching prototype, so the comparison with production systems
is conceptual rather than a performance benchmark.

\begin{table}[H]
\centering
\begin{tabular}{p{0.28\linewidth} p{0.25\linewidth} p{0.34\linewidth}}
\hline
\textbf{MLStore-Lite layer} & \textbf{Production reference} & \textbf{Shared idea} \\
\hline
Storage engine & RocksDB / LevelDB & WAL, memtable, SSTables, compaction \\
Sharding and replication & Cassandra & partitioned key space, replicated data \\
Event log & Kafka & append-only events, producers, consumers, offsets \\
Batch processing & Spark / MapReduce & map, shuffle, reduce over finite data \\
Stream processing & Flink / Kafka Streams & continuous event processing and windowed state \\
Integrated workflow & Databricks / managed platforms & end-to-end data and ML infrastructure \\
Inference extension & Feature stores / model serving systems & online feature retrieval and prediction logging \\
\hline
\end{tabular}
\caption{Conceptual comparison between MLStore-Lite and production systems.}
\end{table}

The main difference is scale and operational complexity. Production tools run
across real machines, handle concurrent users, recover from many failure modes,
and optimize for performance. MLStore-Lite runs locally and focuses on making
the core ideas understandable.

\section{Limitations}

The most important limitation is that the project is local. Nodes are Python
objects and directories, not independent networked processes. There is no real
network transport, automatic leader election, background repair, distributed
query execution, or production-grade monitoring.

The evaluation is also intentionally small. The measurements are useful for
checking that the system runs and for discussing relative behavior, but they are
not benchmarks against real systems. The synthetic workloads are small enough to
keep the project readable and runnable on a laptop.

These limitations are acceptable because the goal is educational. The project
prioritizes clarity of architecture over production completeness.

\section{Conclusion}

MLStore-Lite implements a compact data infrastructure prototype for ML-style
features. The project starts from a single-node storage engine and gradually
adds replication, sharding, batch processing, stream processing, integration,
observability, and online model inference.

The final result shows how raw events can become derived feature values, and how
those values can be stored in a sharded replicated backend, served to a model,
and logged as predictions. It also shows why evaluation and observability
matter: once a system has several layers, it needs logs, measurements, and
reproducible records to be understandable.
